
\documentclass{article}
\usepackage{colortbl}
\usepackage{makecell}
\usepackage{multirow}
\usepackage{supertabular}

\begin{document}

\newcounter{utterance}

\centering \large Interaction Transcript for game `cladder', experiment `full\_v1.5\_default', episode 1468 with Qwen3.5{-}9B.
\vspace{24pt}

{ \footnotesize  \setcounter{utterance}{1}
\setlength{\tabcolsep}{0pt}
\begin{supertabular}{c@{$\;$}|p{.15\linewidth}@{}p{.15\linewidth}p{.15\linewidth}p{.15\linewidth}p{.15\linewidth}p{.15\linewidth}}
    \# & \multicolumn{2}{c}{Player} && \multicolumn{2}{c}{Game Master} \\
    \hline

    \theutterance \stepcounter{utterance}  
    & & & \multicolumn{4}{p{0.6\linewidth}}{
        \cellcolor[rgb]{0.9,0.9,0.9}{
            \makecell[{{p{\linewidth}}}]{
                \texttt{\tiny{[P1$\langle$GM]}}
                \texttt{You are an expert in causal inference. The following question is not a typical commonsense query, but rather a meticulously designed question created by a professor specializing in causal inference, intended to assess the students' mastery of the course content.} \\
\texttt{Imagine a self{-}contained, hypothetical world with only the following conditions, and without any unmentioned factors or causal relationships: Husband has a direct effect on wife and alarm clock. Wife has a direct effect on alarm clock. The overall probability of alarm set by husband is 8\%. For husbands that don't set the alarm, the probability of ringing alarm is 75\%. For husbands that set the alarm, the probability of ringing alarm is 16\%. Is ringing alarm more likely than silent alarm overall?} \\
\texttt{Start your answer with 'Yes' or 'No', followed by additional reasoning or evidence to support your explanation.} \\
            }
        }
    }
     \\ \\

    \theutterance \stepcounter{utterance}  
    & \multicolumn{4}{p{0.6\linewidth}}{
        \cellcolor[rgb]{0.561,0.737,0.561}{
            \makecell[{{p{\linewidth}}}]{
                \texttt{\tiny{[P1$\rangle$GM]}}
                \texttt{No} \\
\\ 
\texttt{**Reasoning:**} \\
\\ 
\texttt{To determine if the alarm is more likely to ring than remain silent overall, we must calculate the total probability of the alarm ringing, \$P(\textbackslash{}text\{Ring\})\$, and compare it to the probability of it being silent, \$P(\textbackslash{}text\{Silent\}) = 1 {-} P(\textbackslash{}text\{Ring\})\$.} \\
\\ 
\texttt{Based on the conditions provided in this hypothetical world:} \\
\texttt{1.  **Husband's Action**: The probability that the husband sets the alarm is given as \$P(\textbackslash{}text\{Set\}) = 0.08\$ (8\%). Consequently, the probability that the husband does *not* set the alarm is \$P(\textbackslash{}text\{No Set\}) = 1 {-} 0.08 = 0.92\$ (92\%).} \\
\texttt{2.  **Conditional Probabilities**:} \\
\texttt{    *   If the husband sets the alarm, the probability of ringing is \$P(\textbackslash{}text\{Ring\} | \textbackslash{}text\{Set\}) = 0.16\$ (16\%).} \\
\texttt{    *   If the husband does not set the alarm, the probability of ringing is \$P(\textbackslash{}text\{Ring\} | \textbackslash{}text\{No Set\}) = 0.75\$ (75\%).} \\
\\ 
\texttt{*Note: The information regarding the "wife" having a direct effect on the alarm clock and the husband having a direct effect on the wife is extraneous to the calculation of the marginal probability of the alarm ringing, as} \\
            }
        }
    }
    & & \\ \\

    \theutterance \stepcounter{utterance}  
    & & & \multicolumn{2}{p{0.3\linewidth}}{
        \cellcolor[rgb]{0.9,0.9,0.9}{
            \makecell[{{p{\linewidth}}}]{
                \texttt{\tiny{[GM$|$GM]}}
                \texttt{No} \\
            }
        }
    }
    & & \\ \\

    \theutterance \stepcounter{utterance}  
    & & & \multicolumn{2}{p{0.3\linewidth}}{
        \cellcolor[rgb]{0.9,0.9,0.9}{
            \makecell[{{p{\linewidth}}}]{
                \texttt{\tiny{[GM$|$GM]}}
                \texttt{yes} \\
            }
        }
    }
    & & \\ \\

    \theutterance \stepcounter{utterance}  
    & & & \multicolumn{2}{p{0.3\linewidth}}{
        \cellcolor[rgb]{0.9,0.9,0.9}{
            \makecell[{{p{\linewidth}}}]{
                \texttt{\tiny{[GM$|$GM]}}
                \texttt{game\_result = LOSE} \\
            }
        }
    }
    & & \\ \\

\end{supertabular}
}

\end{document}
